% !TEX root = 0000.tex
% Anexo C
% Contenido tomado de main_2.tex (líneas 3689-3877) sin cambios de texto.

% El texto está en inglés: reglas de corte de palabra del inglés
\begin{hyphenrules}{english}

\chapter{Anexo: Derivations for the Optimization Problem of PINNs}\chaptermark{Optimization Derivations}\label{app:optimization}

This appendix contains the intermediate computations of Sections~\ref{sec:optimization} and
\ref{sec:gradient-pathologies}.

% ---------------------------------------------------------------------
\section{Gradients and KKT System in Matrix Form}\label{app-sec:kkt-gradients}

For the KKT conditions \eqref{eq:kkt-1}--\eqref{eq:kkt-2}, expanding the gradient of the initial
condition term yields
\begin{align}
  \nabla_\theta\Loss_{ic}(\theta)
  &=\nabla_\theta\Big(\frac{1}{\abs{\T_{ic}}}\sum_{\mathbf{x}\in\T_{ic}}\abs{u_\theta(\mathbf{x},0)-g_0(\mathbf{x})}^2\Big) \label{eq:grad-ic-1}\\
  &=\frac{2}{\abs{\T_{ic}}}\sum_{\mathbf{x}\in\T_{ic}}\big(u_\theta(\mathbf{x},0)-g_0(\mathbf{x})\big)\nabla_\theta u_\theta(\mathbf{x},0). \label{eq:grad-ic-2}
\end{align}
The terms $2u_\theta\nabla_\theta u_\theta$ and $2g_0\nabla_\theta u_\theta$ appear; both are nonlinear in
$\theta$, because $u_\theta$ is nonlinear and so is its derivative. The same behavior is found in the
gradient of the boundary condition. For the physics loss, the gradient is even more complex because of
the differential operator; for the Klein--Gordon operator
$\mathcal{D}_{KG}=\partial_{tt}-\nabla^2+m^2$,
\begin{align}
  \nabla_\theta\Loss_f(\theta)
  &=\frac{1}{\abs{\T_f}}\sum_{(\mathbf{x},t)\in\T_f}\nabla_\theta\abs{\mathcal{D}_{KG}[u_\theta](\mathbf{x},t)}^2 \label{eq:grad-f-1}\\
  &=\frac{2}{\abs{\T_f}}\sum_{(\mathbf{x},t)\in\T_f}\mathcal{D}_{KG}[u_\theta](\mathbf{x},t)\;
    \mathcal{D}_{KG}\big[\nabla_\theta u_\theta\big](\mathbf{x},t). \label{eq:grad-f-2}
\end{align}
Notice that $\nabla_\theta\Loss_f$ requires differentiating with respect to the parameters through the
temporal and spatial derivatives of the network ($\partial_{tt}$, $\nabla^2$), which increases the
computational complexity compared with a standard neural network.

\paragraph{The KKT system in matrix form.}
Writing $N_{ic}=\abs{\T_{ic}}$, $N_{bc}=\abs{\T_{bc}}$, $N_f=\abs{\T_f}$, the expanded KKT conditions
show that the optimizer must solve a system of high-dimensional nonlinear equations (the
dimensionality is induced by the neural network):
\begin{align}
  \frac{2}{N_{ic}}J_{ic}^{T}R_{ic}(\theta)+\frac{2}{N_{bc}}J_{bc}^{T}R_{bc}(\theta)
  +\frac{2v}{N_f}J_f^{T}R_f(\theta)&=0, \label{eq:kkt-matrix-1}\\
  \Loss_f(\theta)&=0, \label{eq:kkt-matrix-2}
\end{align}
where $R_{ic}(\theta)\in\mathbb{R}^{N_{ic}}$ has entries $u_\theta(\mathbf{x}_i,0)-g_0(\mathbf{x}_i)$;
$J_{ic}(\theta)\in\mathbb{R}^{N_{ic}\times P}$ has rows $\nabla_\theta u_\theta(\mathbf{x}_i,0)^{T}$ (analogously for
$R_{bc},J_{bc}$); $R_f(\theta)\in\mathbb{R}^{N_f}$ has entries $\mathcal{D}_{KG}[u_\theta](\mathbf{x}_k,t_k)$; and
$J_f(\theta)\in\mathbb{R}^{N_f\times P}$ has rows $\mathcal{D}_{KG}[\nabla_\theta u_\theta](\mathbf{x}_k,t_k)^{T}$.

\paragraph{One multiplier per collocation point.}
Since the single constraint $\Loss_f(\theta)=0$ is degenerate (Section~\ref{sec:optimization}), the natural
hard-constrained problem imposes $R_f(\mathbf{x}_k,t_k;\theta)=0$ at each of the $N_f$ collocation points,
with $N_f$ multipliers $v_k$, and Equation~\eqref{eq:kkt-matrix-1} becomes
$\frac{2}{N_{ic}}J_{ic}^{T}R_{ic}+\frac{2}{N_{bc}}J_{bc}^{T}R_{bc}+J_f^{T}\mathbf{v}=0$. Whether this
system has a solution depends on whether the network is expressive enough to reach that optimum.

% ---------------------------------------------------------------------
\section{Non-Convexity of the Loss}\label{app-sec:nonconvexity}

For the initial condition term,
\begin{equation}
  \nabla^2_\theta\Loss_{ic}(\theta)=\frac{2}{N_{ic}}\sum_{\mathbf{x}\in\T_{ic}}
  \Big[\nabla_\theta u_\theta(\mathbf{x},0)\,\nabla_\theta u_\theta(\mathbf{x},0)^{T}
  +\big(u_\theta(\mathbf{x},0)-g_0(\mathbf{x})\big)\nabla^2_\theta u_\theta(\mathbf{x},0)\Big].
  \label{eq:hessian-ic}
\end{equation}
The first term, an outer product, is always positive semidefinite (it is the Gauss--Newton part). The
interesting analysis is in the second term, for which we do not know the sign of
$u_\theta(\mathbf{x},0)-g_0(\mathbf{x})$, and $\nabla^2_\theta u_\theta$ is, in general, indefinite. Near the optimum this
factor is small, but globally nothing guarantees positivity. In fact, writing $\Loss_{ic}=F\circ G$ with
$G(\theta)=u_\theta(\cdot,0)-g_0$ and $F=\norm{\cdot}_2^2$, the composition rules of
Boyd and Vandenberghe~\cite{boyd2004} would require $G$ to be affine (or convex with $F$ non-decreasing); but even for a
single layer $u_\theta(\mathbf{y})=b+V\sigma(c+\mathbf{w}\cdot\mathbf{y})$, the dependence on $(\mathbf{w},c)$ passes
through an activation that is neither convex nor concave, so $\Loss_{ic}$ is not convex. The cost function
is non-convex, and non-convex optimization methods must be used.

% ---------------------------------------------------------------------
\section{Infeasible-Start Newton Method}\label{app-sec:newton}

The Newton step $\nabla^2\mathscr{L}\,\delta=-\nabla\mathscr{L}$ applied to both variables $(\theta,v)$
with the linearized constraint is the linear system \cite{boyd2004}
\begin{equation}
  \begin{pmatrix}
    \nabla^2_\theta\Loss_{ic}+\nabla^2_\theta\Loss_{bc}+v\nabla^2_\theta\Loss_f & \nabla_\theta\Loss_f\\
    \nabla_\theta\Loss_f^{T} & 0
  \end{pmatrix}
  \begin{pmatrix}\delta\theta_{\mathrm{nt}}\\ \delta v_{\mathrm{nt}}\end{pmatrix}
  =-\begin{pmatrix}\nabla_\theta\Loss_{ic}+\nabla_\theta\Loss_{bc}+v\nabla_\theta\Loss_f\\ \Loss_f\end{pmatrix}.
  \label{eq:newton-system}
\end{equation}
Its matrix needs the Hessian of the approximating function $u_\theta$ for the losses of the initial and
boundary data, and the Hessian of the physics loss, which in turn contains spatial and temporal
derivatives of $u_\theta$; its cost is estimated in Section~\ref{sec:optimization}.

% ---------------------------------------------------------------------
\section{Stiffness of Gradient Descent}\label{app-sec:stiffness}

For the steepest descent step \eqref{eq:gd}, the Taylor expansion of the loss function around the
current parameters, up to second order, is
\begin{equation}
  \Loss(\theta_{n+1})=\Loss(\theta_n)+(\theta_{n+1}-\theta_n)^{T}\nabla_\theta\Loss(\theta_n)
  +\frac12(\theta_{n+1}-\theta_n)^{T}\nabla^2_\theta\Loss(\xi)(\theta_{n+1}-\theta_n).
  \label{eq:taylor}
\end{equation}
Using the gradient descent update rule $\theta_{n+1}-\theta_n=-\eta\nabla_\theta\Loss(\theta_n)$, the
loss difference becomes
\begin{equation}
  \Loss(\theta_{n+1})-\Loss(\theta_n)=-\eta\norm{\nabla_\theta\Loss(\theta_n)}_2^2
  +\frac12\eta^2\nabla_\theta\Loss(\theta_n)^{T}\nabla^2_\theta\Loss(\xi)\nabla_\theta\Loss(\theta_n).
  \label{eq:loss-diff}
\end{equation}
It is worth analyzing each term. The first one is the squared gradient times the learning rate. For the
second one, let $\mathbf{x}=\nabla_\theta\Loss(\theta_n)/\norm{\nabla_\theta\Loss(\theta_n)}$ be the unit
direction of the gradient. Since the Hessian $\nabla^2_\theta\Loss(\xi)$ is a real symmetric matrix, by the
spectral theorem \cite{trefethen1997} it can be diagonalized by an orthogonal matrix $Q$ that contains
its eigenvectors, $\nabla^2_\theta\Loss(\xi)=Q^{T}\diag(\lambda_1,\dots,\lambda_P)Q$, with
$\lambda_1\le\dots\le\lambda_P$. With $\mathbf{y}=Q\mathbf{x}$,
\begin{equation}
  \nabla_\theta\Loss(\theta_n)^{T}\nabla^2_\theta\Loss(\xi)\nabla_\theta\Loss(\theta_n)
  =\norm{\nabla_\theta\Loss(\theta_n)}^2\,\mathbf{y}^{T}\diag(\lambda_1,\dots,\lambda_P)\mathbf{y}
  =\norm{\nabla_\theta\Loss(\theta_n)}^2\sum_{i=1}^{P}\lambda_iy_i^2.
  \label{eq:hessian-quadratic}
\end{equation}
Substituting in Equation~\eqref{eq:loss-diff}, we obtain Equation~\eqref{eq:stiffness}.

% ---------------------------------------------------------------------
\section{Bounds on the Gradients}\label{app-sec:grad-bounds}

Wang et al.~\cite{wang2021gradient} develop the analysis of bounds starting from inequalities of Sobolev type
for $\nabla_\theta\Loss_f$ and $\nabla_\theta\Loss_{\mathrm{data}}$:
\begin{equation}
  \norm{\nabla_\theta\Loss_f(\theta)}_{L^\infty}\le C\norm{\nabla^2_\theta\Loss_f(\theta)}_{L^\infty},\qquad
  \norm{\nabla_\theta\Loss_{\mathrm{data}}(\theta)}_{L^\infty}\le C\norm{\nabla^2_\theta\Loss_{\mathrm{data}}(\theta)}_{L^\infty}.
  \label{eq:grad-bounds}
\end{equation}
\verificar{exact form and hypotheses of these inequalities in Wang et al.\ (2021)}
This imposes that the maximum value of the gradient is bounded by the maximum value of the curvature of
the loss function. To connect this with the eigenvalues of the Hessian, we use the equivalence of matrix
norms \cite{trefethen1997}: for $H\in\mathbb{R}^{P\times P}$,
$P^{-1/2}\norm{H}_\infty\le\norm{H}_2\le P^{1/2}\norm{H}_\infty$. Then the inequalities are directly related
to the eigenvalues of the Hessians of the physics and data losses:
\begin{align}
  \norm{\nabla_\theta\Loss_f(\theta)}_{L^\infty}&\le C\sqrt{P}\,\norm{\nabla^2_\theta\Loss_f(\theta)}_{2}
  =C\sqrt{P}\,\sqrt{\lambda_{\max}(H_f^{T}H_f)}, \label{eq:bound-f}\\
  \norm{\nabla_\theta\Loss_{\mathrm{data}}(\theta)}_{L^\infty}&\le C\sqrt{P}\,\norm{\nabla^2_\theta\Loss_{\mathrm{data}}(\theta)}_{2}
  =C\sqrt{P}\,\sqrt{\lambda_{\max}(H_{\mathrm{data}}^{T}H_{\mathrm{data}})}. \label{eq:bound-data}
\end{align}
So the size of the gradient is bounded by the largest eigenvalue of its curvature. A large eigenvalue of
$H^{T}H$ allows large values of the gradient, while a small eigenvalue restricts the gradient to small
values. Therefore, the imbalance between gradients can be reaffirmed.


% ---------------------------------------------------------------------
\section{Fisher Matrix and KL Divergence}\label{app:fisher}

Let the Kullback--Leibler (KL) divergence between $P(\cdot\mid\theta)$ and $P(\cdot\mid\theta+\delta)$ be
\cite{pascanu2013}
\[
  \mathrm{KL}(\theta\,\|\,\theta+\delta)=\mathbb{E}_{\mathbf{x}\sim P(\cdot\mid\theta)}\big[\log P(\mathbf{x}\mid\theta)-\log P(\mathbf{x}\mid\theta+\delta)\big].
\]
Expanding $\log P(\mathbf{x}\mid\theta+\delta)$ to second order, we obtain
\[
  \log P(\mathbf{x}\mid\theta+\delta)\approx\log P(\mathbf{x}\mid\theta)+\nabla_\theta\log P(\mathbf{x}\mid\theta)^{T}\delta
  +\frac12\delta^{T}\nabla^2_\theta\log P(\mathbf{x}\mid\theta)\,\delta,
\]
and substituting,
\begin{equation}
  \mathrm{KL}(\theta\,\|\,\theta+\delta)\approx-\mathbb{E}\big[\nabla_\theta\log P(\mathbf{x}\mid\theta)\big]^{T}\delta
  -\frac12\delta^{T}\,\mathbb{E}\big[\nabla^2_\theta\log P(\mathbf{x}\mid\theta)\big]\,\delta.
  \label{eq:app-kl}
\end{equation}
The first term vanishes for every $\theta$ (not only at the optimum), because of the following properties:
\[
  \int P(\mathbf{x}\mid\theta)\,d\mathbf{x}=1\ \Rightarrow\ \int\nabla_\theta P(\mathbf{x}\mid\theta)\,d\mathbf{x}=0,
\]
\[
  \nabla_\theta\log P=\frac{\nabla_\theta P}{P}\ \Rightarrow\
  \mathbb{E}\big[\nabla_\theta\log P(\mathbf{x}\mid\theta)\big]=\int P\,\nabla_\theta\log P\,d\mathbf{x}=0.
\]
Differentiating the last identity once more with respect to $\theta$,
\[
  0=\int\Big(\nabla_\theta P\,\nabla_\theta\log P^{T}+P\,\nabla^2_\theta\log P\Big)\,d\mathbf{x}
   =\int P\Big(\nabla_\theta\log P\,\nabla_\theta\log P^{T}+\nabla^2_\theta\log P\Big)\,d\mathbf{x},
\]
that is, the expectation of the second moment satisfies
\[
  \mathbb{E}\big[\nabla_\theta\log P\,\nabla_\theta\log P^{T}\big]=-\mathbb{E}\big[\nabla^2_\theta\log P\big].
\]
Therefore,
\[
  \mathrm{KL}(\theta\,\|\,\theta+\delta)\approx\frac12\delta^{T}\,
  \mathbb{E}\big[\nabla_\theta\log P(\mathbf{x}\mid\theta)\,\nabla_\theta\log P(\mathbf{x}\mid\theta)^{T}\big]\,\delta
  =\frac12\delta^{T}F\delta.
\]

\end{hyphenrules}
